\documentclass[11pt,a4paper]{article}
\usepackage[utf8]{inputenc}
\usepackage[T1]{fontenc}
\usepackage[spanish]{babel}
\usepackage{geometry}
\geometry{margin=2.5cm}
\usepackage{graphicx}
\usepackage{booktabs}
\usepackage{listings}
\usepackage{xcolor}
\usepackage{hyperref}
\usepackage{enumitem}

\lstset{
    basicstyle=\ttfamily\footnotesize,
    backgroundcolor=\color{gray!10},
    breaklines=true,
    frame=single,
    showstringspaces=false,
    numbers=left,
    numberstyle=\tiny\color{gray},
    captionpos=b
}

\title{\textbf{Proyecto 1 — Base de Datos II}\\[0.4em]
\large Organización de archivos: índices SQL sobre B\textsuperscript{+}-Tree, Sequential File, Extendible Hash y R-Tree}
\author{Grupo BD2 — UTEC 26-1\\\small{[completar integrantes]}}
\date{\today}

\begin{document}
\maketitle

\begin{abstract}
Se diseñó e implementó un mini-DBMS que expone un dialecto SQL reducido
sobre cuatro estructuras de archivo: \textbf{B\textsuperscript{+}-Tree},
\textbf{Sequential File}, \textbf{Extendible Hash} y \textbf{R-Tree
espacial}. Cada técnica se evalúa midiendo accesos a disco
(\emph{disk reads/writes}) y tiempo de ejecución para inserción, búsqueda
puntual, búsqueda por rango y --- en el caso del R-Tree --- búsqueda por
radio y $k$-vecinos más cercanos. El R-Tree adicionalmente integra control
de concurrencia 2PL con detección de \emph{deadlocks} y \emph{Write-Ahead
Log} para recuperación ante fallos.
\end{abstract}

\section{Introducción}

\subsection{Objetivos}
\begin{itemize}[nosep]
    \item Implementar cuatro técnicas clásicas de indexación en disco con
    operaciones \texttt{search}, \texttt{add}, \texttt{remove} y, donde
    aplique, \texttt{range\_search}.
    \item Construir un parser SQL minimalista y un ejecutor que despache la
    operación a la estructura adecuada según el índice declarado en
    \texttt{CREATE TABLE}.
    \item Comparar empíricamente el costo de I/O y el tiempo de respuesta
    entre las cuatro estructuras usando un dataset de libros.
    \item En la rama del R-Tree: incorporar 2PL con \emph{lock coupling},
    detección de \emph{deadlocks} por wait-for graph, y \emph{Write-Ahead
    Log} ARIES-simplificado.
\end{itemize}

\subsection{Dominio de datos}
Se utiliza el dataset \texttt{books.csv} con los campos
\texttt{book\_key, title, author, pages, average\_rating,
published\_date}. Para consultas espaciales se usan coordenadas
\texttt{(longitud, latitud)} con TIDs sintéticos.

\section{Arquitectura del sistema}

\begin{verbatim}
            +---------------------+
            |  Cliente / Frontend |
            +----------+----------+
                       |
                       v
            +---------------------+
            |     parser_sql.py   |  Scanner + Parser → AST
            +----------+----------+
                       |
                       v
            +---------------------+
            |     ejecutor.py     |  Catálogo de tablas + dispatch
            +-+-------+--------+--+
              |       |        |   \
              v       v        v    v
         BPlusTree SeqFile ExtHash RTree (+WAL +LockManager)
\end{verbatim}

Todos los índices comparten una interfaz informal:
\texttt{search(k)}, \texttt{add(rec)}/\texttt{insert(rec)}, \texttt{remove(k)},
\texttt{range\_search(...)}, contadores de \texttt{disk\_reads/writes}.

\section{Estructuras implementadas}

\subsection{B\textsuperscript{+}-Tree}
Páginas de tamaño fijo. Las hojas almacenan los registros completos y se
enlazan en lista simple para soportar \texttt{range\_search} eficiente.
Los nodos internos almacenan claves separadoras y punteros a página.
Inserciones causan \emph{splits} con propagación a la raíz; eliminaciones
implementadas con \emph{borrado lógico} (sin fusión inmediata).

\textbf{Costos esperados}: $O(\log_B n)$ para búsqueda puntual,
$O(\log_B n + k/B)$ para rango.

\subsection{Sequential File}
Archivo principal ordenado más archivo auxiliar para overflows. La
búsqueda usa búsqueda binaria sobre el archivo principal y luego un
recorrido secuencial del auxiliar. Periódicamente se hace \texttt{rebuild}
para reabsorber el auxiliar.

\textbf{Trade-off}: lecturas baratas por localidad pero inserciones
costosas; \texttt{k\_desorted} controla la frecuencia de rebuild.

\subsection{Extendible Hash}
Directorio en memoria de tamaño $2^{D}$ donde $D$ es la profundidad
global. Cada bucket en disco con \texttt{local\_depth}. Cuando un bucket
se llena se hace \emph{split} con redistribución por bit relevante; si
\texttt{local\_depth = global\_depth} se duplica el directorio.

\textbf{Limitación}: no soporta \texttt{range\_search} (consecuencia
inherente del hashing).

\textbf{Bugs corregidos durante la implementación}:
\begin{enumerate}[nosep]
    \item Variable local llamada \texttt{bit\_valor} sobreescribía la función homónima en el scope del split.
    \item En el re-puntero del directorio se aplicaba \texttt{bit\_valor(i, new\_ld)} sobre el índice $i$, cuando $i$ ya \emph{es} el prefijo de bits y solo había que extraer el bit relevante.
    \item \texttt{insert} llamaba a \texttt{\_split} una sola vez sin reintento: si dos claves caían en el mismo bucket post-split (colisión total de bits) la inserción se perdía. Reemplazado por bucle con tope de reintentos.
\end{enumerate}

\subsection{R-Tree}
Árbol balanceado sobre MBRs (Minimum Bounding Rectangles).
Implementación de splits cuadrática (algoritmo de Guttman) y descenso por
\emph{minimum enlargement}. Soporta:
\begin{itemize}[nosep]
    \item \texttt{range\_search(x, y, r)}: poda por intersección de
    círculo con MBR.
    \item \texttt{knn(x, y, k)}: \emph{best-first search} con cola de
    prioridad sobre distancias mínimas a MBRs.
    \item \texttt{insert(lon, lat, tid)}: TID es par
    $(\text{page\_id}, \text{slot\_id})$ que apunta al registro real.
\end{itemize}

\section{Concurrencia y recuperación (R-Tree)}

\subsection{Two-Phase Locking con lock coupling}
\textsf{LockManager} a nivel de página con modos S (compartido) y X
(exclusivo). El descenso por el árbol usa \emph{lock coupling}
(\textit{crabbing}): se mantiene un S-lock sobre el padre solo hasta
asegurar el lock del hijo.

\subsection{Detección de \emph{deadlocks}}
Wait-for graph que se reconstruye en cada espera. DFS detecta ciclos; la
víctima es la transacción con menor \texttt{tid}, que es marcada
(\texttt{aborted = True}) y libera todos sus locks. Las demás
transacciones la observan vía \texttt{Condition.wait(timeout=0.05)} y
proceden.

\subsection{Write-Ahead Log}
Esquema ARIES simplificado. Antes de cada modificación se escribe
\texttt{(tid, pid, before\_image, after\_image)} al log. Recovery aplica
REDO de transacciones \emph{committed} y UNDO de las que no llegaron a
\texttt{COMMIT}.

\section{Parser y ejecutor SQL}

El subset SQL soportado:
\begin{lstlisting}[language=SQL,caption={Gramática soportada}]
CREATE TABLE t (col TIPO [INDEX TECNICA], ...) [FROM FILE "x.csv"];
SELECT * FROM t WHERE col = v;
SELECT * FROM t WHERE col BETWEEN a AND b;
SELECT * FROM t WHERE col IN (POINT(x,y), RADIUS r);
SELECT * FROM t WHERE col IN (POINT(x,y), K n);
INSERT INTO t VALUES (v1, v2, ...);
DELETE FROM t WHERE col = v;
\end{lstlisting}

Las técnicas válidas son \texttt{BPTREE}, \texttt{SEQUENTIAL},
\texttt{HASH}, \texttt{RTREE}. El parser produce un AST de
\textit{dataclasses} (\texttt{NodoCreateTable},
\texttt{NodoSelectPuntual}, etc.) que el \textsf{Ejecutor} despacha al
índice correspondiente y devuelve un \texttt{ResultadoEjecucion} con
filas, \texttt{disk\_reads/writes} y tiempo en milisegundos.

\section{Pruebas}

\subsection{Suite de tests automatizados}
\begin{itemize}[nosep]
    \item 71 tests unitarios sobre el R-Tree
    (\texttt{test\_rtree\_ops}, \texttt{test\_wal},
    \texttt{test\_lock\_manager}, \texttt{test\_recovery},
    \texttt{test\_concurrency}, \dots).
    \item 3 tests de integración del ejecutor SQL contra B\textsuperscript{+}-Tree, Hash y R-Tree.
    \item Total: \textbf{74 tests, 100\% pasando} en $\sim 5\,$s.
\end{itemize}

\subsection{Demo de concurrencia}
\texttt{examples/concurrent\_demo.py} lanza dos threads que insertan y
consultan en rangos disjuntos del mismo árbol. El \texttt{event\_cb} del
LockManager registra cada \texttt{LOCK}/\texttt{WAIT} con \emph{timestamp}
en \texttt{concurrent\_demo.log}, mostrando la contención real, los
\texttt{COMMIT} por lote y los \texttt{ABORT} por \emph{deadlock}.

\section{Resultados experimentales (preliminares)}

\begin{table}[h]
\centering
\small
\begin{tabular}{lrrrr}
\toprule
\textbf{Operación} & \textbf{B\textsuperscript{+}-Tree} & \textbf{SeqFile} & \textbf{Hash} & \textbf{R-Tree} \\
\midrule
\texttt{search} puntual (I/O) & 3 & $O(\log n)$ + aux & 1--2 & --- \\
\texttt{range\_search} (I/O)  & 3 + hojas & lineal en rango & N/A & poda MBR \\
\texttt{insert} (I/O)         & 6 & 1 (al aux) & 1--2 & $O(\log n)$ \\
\bottomrule
\end{tabular}
\caption{Costos observados en demos pequeñas (n=10). Los datos finales se obtendrán con \texttt{books.csv} completo.}
\end{table}

\textbf{Pendiente}: ejecutar benchmark sobre $n \in \{10^3, 10^4, 10^5\}$
con \texttt{books.csv} y graficar tiempo y \texttt{total\_io} vs $n$.

\section{Conclusiones (preliminares)}

\begin{itemize}[nosep]
    \item B\textsuperscript{+}-Tree es la opción más balanceada: soporta puntual y rango con costo logarítmico predecible.
    \item Sequential File es competitivo para cargas mayoritariamente de lectura, pero requiere \emph{rebuild} periódico.
    \item Extendible Hash es óptimo para búsqueda puntual (1--2 I/O) pero no admite rango.
    \item R-Tree es la única estructura del proyecto que soporta consultas espaciales (radio y $k$-NN).
    \item La integración del WAL y el LockManager en el R-Tree no degradó las pruebas single-threaded existentes y permitió 2 threads concurrentes con detección automática de \emph{deadlocks}.
\end{itemize}

\section{Trabajo pendiente}
\begin{itemize}[nosep]
    \item Benchmark formal con \texttt{books.csv} completo y gráficas comparativas.
    \item Frontend web (rama \texttt{frontend}) que consuma el ejecutor.
    \item Integración del parser desde el frontend (POST de SQL → tabla de resultados con métricas).
    \item Video de demostración.
\end{itemize}

\appendix
\section{Estructura del repositorio}
\begin{verbatim}
BD2-Project/
├── parser_sql.py          ← scanner + parser SQL
├── ejecutor.py            ← dispatcher AST → índice
├── bplustree.py           ← B+-Tree
├── sequential_file.py     ← Sequential File con archivo aux
├── extendible_hash.py     ← Extendible Hash
├── rtree/
│   ├── rtree.py           ← R-Tree con lock coupling
│   ├── page_store.py      ← gestor de páginas + WAL
│   ├── lock_manager.py    ← 2PL + deadlock detection
│   ├── wal.py             ← Write-Ahead Log ARIES-simplificado
│   ├── transaction.py     ← context manager begin/commit/abort
│   └── geometry.py        ← MBR, TID, primitivas geométricas
├── examples/
│   └── concurrent_demo.py ← demo de 2 threads con log timestamped
└── tests/                 ← 74 tests
\end{verbatim}

\end{document}
